%!TEX root = main.tex

%In the sequel, we first introduce a concept of expression trees to specify the correctness property of {\tt expand\_branch}, then prove that {\tt expand\_branch} satisfies the property. 

\paragraph*{Specification}
As in the other phases, the correctness of {\sf expandWhens} is specified by the formula:
%
$$(tm,\,cm) \models \, {\sf expandWhens}(M) \Rightarrow (tm,\,cm) \models \,M.$$
%
This formula states that if the pair of the typing mapping $tm$ and the connection mapping $ct$ conform to the program produced by {\sf expandWhens} $(M)$,
then the same pair of $tm$ and $ct$ should also conform to the original program $M$.
Note that this requires that an incorrect program is not translated to a correct program;
for example, if $M$ contains scoping errors like in Figure~\ref{fig:example},
{\sf expandWhens} cannot produce an output program.
Therefore, {\sf expandBranch} also checks for such scoping errors.

The correctness proof uses in a mutual induction over statements and statement sequences.
If a statement $s$ and a statement sequence $ss$ satisfy the correctness property, then also the statement sequence $s ; ss$.
Moreover, a statement \prettyll{when} $c$ \prettyll{:} $sst$ \prettyll{else} $ssf$
satisfies the correctness property if both $sst$ and $ssf$ satisfy it.
The mutual induction also requires us to define the two Coq datatypes for statements and statement sequences mutually inductive;
we cannot just use the sequence or list types from the standard Coq library.

The proof compares in detail the constructions and checks executed by {\sf stmtsTmap},
the construction of a connection mapping, and {\sf expandBranch}.
The main maps that should \emph{fit together} are:
$scope$, $cm$ (constructed by {\sf expandBranch}), $cm$ (the connection mapping constructed earlier), $tm$ and $tmap$.
Generally speaking, the domain of every map in this list is a subdomain of the next map
(and the information in the maps needs to be consistent with each other, e.g. sources cannot be connected to).
In particular, the $cm$ constructed by {\sf expandBranch} needs to contain all connections that are in scope
so that {\sf combineBranches} at the end of a subsequent \prettyll{when} statement can choose correctly
even if only one branch assigns a value.
Then, the other branch retains the old value (like \prettyll{posv} in Figure~\ref{fig-lofirrtl-exmp3}).

The main inductive statement of the proof is:
\begin{lemma}
If:
\begin{itemize}
\item   $tmap$ has fully inferred ground types;
\item   the final $tm$ fits together with the final $tmap$;
\item   statement sequence $ss$ only references fully inferred ground types (e.g. it does not contain uninferred constants);
\item   the five maps mentioned fit together;
\item   {\sf stmtTmap}, the construction of the connection mapping, and {sf expandBranch} can be applied to statement sequence {\sf expandBranch}$(ss)$;
\end{itemize}
Then:
\begin{itemize}
\item   an extensionally equal $tm$ and $cm$ are constructed by running the original statement sequence $ss$;
\item   the resulting maps fit together similarly;
\item   the connection map $cm$ constructed by {\sf expandBranch} contains the same additions for $ss$
        as the connection mapping constructed earlier.
\end{itemize}
\end{lemma}

The complexity of the correctness proof lies mostly in the number of cases that need to be distinguished
and in the multitude of bookkeeping information that needs to be dragged through.
